\section{Proof of Section: Embedding Natural Deduction with Refinement Types}

\begin{lemma}[Validity]\label{lemma:validity}
If there exists $e\in\embed{\formula}$\ then \formula is valid.
\end{lemma}
\begin{proof}
We prove the lemma by case analysis in the shape of \formula. 
\begin{itemize}
\item $\formula \equiv \proofterm{p}$.
Since the set
 $\embed{\proofterm{p}} = \{e | \evalsto{p}{\etrue} \}$
 is not empty, then  \evalsto{p}{\etrue}.
% 
\item $\formula \equiv {\formula_1} \rightarrow {\formula_2}$. 
By assumption, there exists an expressions $f$ so that 
$\forall e_x \in \embed{\formula_1}, f\ e_x \in \embed{\formula_2}$. 
%
So, if there exists an expression 
$e_1 \in \embed{\formula_1}$ that makes $\formula_1$ valid 
then $f\ e_1$ makes $\formula_2$ valid. 
\item $\formula \equiv {\formula} \rightarrow {\efalse}$. 
By assumption, there exists an expressions $f$ so that 
$\forall e_x \in \embed{\formula}, f\ e_x \in \embed{\proofterm{\efalse}}$. 
%
So, if there exists an expression 
$e_1 \in \embed{\formula}$ that makes $\formula$ valid 
then $f\ e_1$ makes $\proofterm{\efalse}$ valid, which is impossible, 
thus \formula cannot be valid. 
\item $\formula \equiv \tyand{\formula_1}{\formula_2}$. 
If there exists a total expression $e \in \embed{\formula}$
then $e$ evaluates to \exprand{e_1}{e_2}. 
\item $\formula \equiv \tyor{\formula_1}{\formula_2}$. 
If there exists a total expression $e \in \embed{\formula}$
then $e$ evaluates to either 
\exprorleft{e'} or \exprorright{e'}. 
\item $\formula \equiv \tyall{x}{\typ}{\formula}$.
By assumption, there exists an expressions $f$ so that 
$\forall e_x \in \embed{\typ}, f\ e_x \in \embed{\formula\subst{x}{e_x}}$. 
%
So, if there exists an expression 
$e_1 \in \embed{\typ}$
then $f\ e_1$ makes $\formula\subst{x}{e_1}$ valid. 
\item $\formula \equiv \tyex{x}{\typ}{\formula}$.
By assumption, there exists an expressions $p$ 
that evaluates to a pair $(e_x, e_y)$ 
so that $e_x\in \embed{\typ}$ and $e_y\in\embed{\formula\subst{x}{e_x}}$.
\end{itemize} 
\end{proof}


\begin{theorem}[Validity]
If \hastype{\emptyset}{\emptyset}{e}{\formula} then \formula is valid.
\end{theorem}
\begin{proof}
By direct implication of Lemma~\ref{lemma:validity} and soundness of \corelan (Theorem~\ref{thm:safety}).
\end{proof}